
\documentclass{article}
\usepackage{colortbl}
\usepackage{makecell}
\usepackage{multirow}
\usepackage{supertabular}

\begin{document}

\newcounter{utterance}

\centering \large Interaction Transcript for game `wordle', experiment `high\_frequency\_words\_no\_clue\_no\_critic', episode 4 with claude{-}opus{-}4{-}8{-}azure.
\vspace{24pt}

{ \footnotesize  \setcounter{utterance}{1}
\setlength{\tabcolsep}{0pt}
\begin{supertabular}{c@{$\;$}|p{.15\linewidth}@{}p{.15\linewidth}p{.15\linewidth}p{.15\linewidth}p{.15\linewidth}p{.15\linewidth}}
    \# & \multicolumn{2}{c}{Player} && \multicolumn{2}{c}{Game Master} \\
    \hline

    \theutterance \stepcounter{utterance}  
    & & & \multicolumn{4}{p{0.6\linewidth}}{
        \cellcolor[rgb]{0.9,0.9,0.9}{
            \makecell[{{p{\linewidth}}}]{
                \texttt{\tiny{[P1$\langle$GM]}}
                \texttt{Você é um mago da linguagem que gosta de adivinhar palavras usando as regras fornecidas.} \\
\\ 
\texttt{Bem{-}vindo ao Wordle! Você tem seis tentativas para adivinhar a palavra{-}alvo, uma palavra válida em português de cinco letras minúsculas. Por favor, use as tags "explicação:" e "palpite:" para fornecer uma explicação concisa para cada palpite.} \\
\\ 
\texttt{Por exemplo, se o seu palpite for "tempo", sua resposta deve ser} \\
\texttt{explicação: esta é uma palavra comum em português de cinco letras, e estou começando meu palpite com esta palavra.} \\
\texttt{palpite: tempo} \\
\\ 
\texttt{Após cada palpite, sua resposta será validada, e você receberá feedback indicando quais letras estão corretas (<verde>), quais letras estão corretas mas na posição errada (<amarelo>) e quais letras estão incorretas (<vermelho>). Esse feedback pode ser útil para determinar quais letras incluir ou excluir no seu próximo palpite.} \\
\\ 
\texttt{Por exemplo, o feedback para "tempo" pode ser:} \\
\texttt{feedback\_palpite: t<amarelo> e<amarelo> m<verde> p<amarelo> o<vermelho>} \\
\\ 
\texttt{A explicação deve conter detalhes sobre como o feedback\_palpite é usado para chegar a um novo palpite.} \\
\\ 
\texttt{Não gere nenhum outro texto.} \\
\\ 
\texttt{Vamos começar com o seu primeiro palpite.} \\
            }
        }
    }
     \\ \\

    \theutterance \stepcounter{utterance}  
    & \multicolumn{4}{p{0.6\linewidth}}{
        \cellcolor[rgb]{0.561,0.737,0.561}{
            \makecell[{{p{\linewidth}}}]{
                \texttt{\tiny{[P1$\rangle$GM]}}
                \texttt{explicação: começo com uma palavra comum que contém muitas vogais e consoantes frequentes para obter informações iniciais.} \\
\texttt{palpite: carne} \\
            }
        }
    }
    & & \\ \\

    \theutterance \stepcounter{utterance}  
    & & & \multicolumn{2}{p{0.3\linewidth}}{
        \cellcolor[rgb]{0.9,0.9,0.9}{
            \makecell[{{p{\linewidth}}}]{
                \texttt{\tiny{[GM$|$GM]}}
                \texttt{attempts = 1} \\
\texttt{target\_word = crime} \\
\texttt{guess = carne} \\
\texttt{guess\_feedback = c<green> a<red> r<yellow> n<red> e<green>} \\
            }
        }
    }
    & & \\ \\

    \theutterance \stepcounter{utterance}  
    & & & \multicolumn{4}{p{0.6\linewidth}}{
        \cellcolor[rgb]{0.9,0.9,0.9}{
            \makecell[{{p{\linewidth}}}]{
                \texttt{\tiny{[P1$\langle$GM]}}
                \texttt{feedback\_palpite: c<verde> a<vermelho> r<amarelo> n<vermelho> e<verde>} \\
\\ 
\texttt{Forneça a sua resposta apenas neste formato.} \\
\texttt{explicação: detalhes} \\
\texttt{palpite: palavra} \\
            }
        }
    }
     \\ \\

    \theutterance \stepcounter{utterance}  
    & \multicolumn{4}{p{0.6\linewidth}}{
        \cellcolor[rgb]{0.561,0.737,0.561}{
            \makecell[{{p{\linewidth}}}]{
                \texttt{\tiny{[P1$\rangle$GM]}}
                \texttt{explicação: C e E estão corretas nas posições 1 e 5. R existe mas não na posição 3. Excluo A e N. Tento "cobre" com R na posição 4.} \\
\texttt{palpite: cobre} \\
            }
        }
    }
    & & \\ \\

    \theutterance \stepcounter{utterance}  
    & & & \multicolumn{2}{p{0.3\linewidth}}{
        \cellcolor[rgb]{0.9,0.9,0.9}{
            \makecell[{{p{\linewidth}}}]{
                \texttt{\tiny{[GM$|$GM]}}
                \texttt{attempts = 2} \\
\texttt{target\_word = crime} \\
\texttt{guess = cobre} \\
\texttt{guess\_feedback = c<green> o<red> b<red> r<yellow> e<green>} \\
            }
        }
    }
    & & \\ \\

    \theutterance \stepcounter{utterance}  
    & & & \multicolumn{4}{p{0.6\linewidth}}{
        \cellcolor[rgb]{0.9,0.9,0.9}{
            \makecell[{{p{\linewidth}}}]{
                \texttt{\tiny{[P1$\langle$GM]}}
                \texttt{feedback\_palpite: c<verde> o<vermelho> b<vermelho> r<amarelo> e<verde>} \\
\\ 
\texttt{Forneça a sua resposta apenas neste formato.} \\
\texttt{explicação: detalhes} \\
\texttt{palpite: palavra} \\
            }
        }
    }
     \\ \\

    \theutterance \stepcounter{utterance}  
    & \multicolumn{4}{p{0.6\linewidth}}{
        \cellcolor[rgb]{0.561,0.737,0.561}{
            \makecell[{{p{\linewidth}}}]{
                \texttt{\tiny{[P1$\rangle$GM]}}
                \texttt{explicação: C e E fixos nas posições 1 e 5. R não pode estar nas posições 3 ou 4, então deve estar na posição 2. Excluo O, B, A, N. Tento "creme".} \\
\texttt{palpite: creme} \\
            }
        }
    }
    & & \\ \\

    \theutterance \stepcounter{utterance}  
    & & & \multicolumn{2}{p{0.3\linewidth}}{
        \cellcolor[rgb]{0.9,0.9,0.9}{
            \makecell[{{p{\linewidth}}}]{
                \texttt{\tiny{[GM$|$GM]}}
                \texttt{attempts = 3} \\
\texttt{target\_word = crime} \\
\texttt{guess = creme} \\
\texttt{guess\_feedback = c<green> r<green> e<red> m<green> e<green>} \\
            }
        }
    }
    & & \\ \\

    \theutterance \stepcounter{utterance}  
    & & & \multicolumn{4}{p{0.6\linewidth}}{
        \cellcolor[rgb]{0.9,0.9,0.9}{
            \makecell[{{p{\linewidth}}}]{
                \texttt{\tiny{[P1$\langle$GM]}}
                \texttt{feedback\_palpite: c<verde> r<verde> e<vermelho> m<verde> e<verde>} \\
\\ 
\texttt{Forneça a sua resposta apenas neste formato.} \\
\texttt{explicação: detalhes} \\
\texttt{palpite: palavra} \\
            }
        }
    }
     \\ \\

    \theutterance \stepcounter{utterance}  
    & \multicolumn{4}{p{0.6\linewidth}}{
        \cellcolor[rgb]{0.561,0.737,0.561}{
            \makecell[{{p{\linewidth}}}]{
                \texttt{\tiny{[P1$\rangle$GM]}}
                \texttt{explicação: Temos C R \_ M E confirmados. A posição 3 precisa de uma vogal diferente de E, A e O. Tento I formando "crime".} \\
\texttt{palpite: crime} \\
            }
        }
    }
    & & \\ \\

    \theutterance \stepcounter{utterance}  
    & & & \multicolumn{2}{p{0.3\linewidth}}{
        \cellcolor[rgb]{0.9,0.9,0.9}{
            \makecell[{{p{\linewidth}}}]{
                \texttt{\tiny{[GM$|$GM]}}
                \texttt{attempts = 4} \\
\texttt{target\_word = crime} \\
\texttt{guess = crime} \\
\texttt{guess\_feedback = c<green> r<green> i<green> m<green> e<green>} \\
            }
        }
    }
    & & \\ \\

    \theutterance \stepcounter{utterance}  
    & & & \multicolumn{2}{p{0.3\linewidth}}{
        \cellcolor[rgb]{0.9,0.9,0.9}{
            \makecell[{{p{\linewidth}}}]{
                \texttt{\tiny{[GM$|$GM]}}
                \texttt{game\_result = WIN} \\
            }
        }
    }
    & & \\ \\

\end{supertabular}
}

\end{document}
